\chapter{Chapter06 — Regularity Comparisons}
%%% generated by the blueprint pipeline from FABL.Chapter06.Pseudorandomness.RegularityComparisons
\section{Overview}

Declarations of the Lean module \texttt{FABL.Chapter06.Pseudorandomness.RegularityComparisons}.

\section{Declarations}

\begin{theorem}[Fourier coefficient bounded by influence]
\lean{FABL.sq_fourierCoeff_le_influence_of_mem}
\proofsource{Analysis-of-Boolean-Functions-by-Ryan-ODonnell}{pdf-ffbe9c7b3306-p048-b001}
\label{FABL.sq_fourierCoeff_le_influence_of_mem}
\leanok
\uses{FABL.fourierCoeff, FABL.influence, FABL.influence_eq_sum_sq_fourierCoeff}
\difficulty{3}
Let $n$ be a nonnegative integer, let $f:\\{-1,1\\}^n\\to\\mathbb R$, and let $S$ be a
set of coordinates containing $i$. Then the square of the Fourier coefficient of $f$ at
$S$ is at most the influence of coordinate $i$ on $f$.
\end{theorem}

\begin{theorem}[Fourier regularity from small influences]
\lean{FABL.HasSmallInfluences.isFourierRegular_sqrt}
\label{FABL.HasSmallInfluences.isFourierRegular_sqrt}
\leanok
\uses{FABL.HasSmallInfluences, FABL.IsFourierRegular, FABL.hasSmallInfluences_iff, FABL.sq_fourierCoeff_le_influence_of_mem}
\difficulty{4}
Let $f:\{-1,1\}^n\to\mathbb R$ and $\epsilon\in\mathbb R$. If every coordinate has
stable influence at correlation parameter $1$ at most $\epsilon$, then every nonempty
set of coordinates $S\subseteq\{1,\ldots,n\}$ satisfies $|\widehat
f(S)|\le\sqrt{\epsilon}$.
\end{theorem}

\begin{theorem}[Regularity of the Mod-Two Inner Product]
\lean{FABL.isFourierRegular_innerProductModTwoBoolean}
\label{FABL.isFourierRegular_innerProductModTwoBoolean}
\leanok
\uses{FABL.BooleanFunction.toReal, FABL.IsFourierRegular, FABL.abs_fourierCoeff_innerProductModTwoBoolean, FABL.innerProductModTwoBoolean}
\difficulty{1}
For every natural number $m$, define a Boolean function on $\{\pm1\}^{2m}$ by decoding
the coordinates into two vectors $a,b\in\mathbb F_2^m$ and returning the indicator of
$\sum_{i=1}^m a_i b_i=1$ in $\mathbb F_2$. Every Fourier coefficient indexed by a
nonempty coordinate set has absolute value at most $2^{-m}$.
\end{theorem}

\begin{theorem}[Influence of the binary inner-product function]
\lean{FABL.influence_innerProductModTwoBoolean_eq_half}
\label{FABL.influence_innerProductModTwoBoolean_eq_half}
\leanok
\uses{FABL.BooleanFunction.toReal, FABL.abs_fourierCoeff_innerProductModTwoBoolean, FABL.fourierCoeff, FABL.influence, FABL.influence_eq_sum_sq_fourierCoeff, FABL.innerProductModTwoBoolean, FABL.sum_pow_card_sub_one_filter_mem}
\difficulty{5}
For every positive integer $m$, define $F:\{-1,1\}^{2m}\to\{-1,1\}$ by splitting the
input into sign vectors corresponding to $u,v\in\mathbb{F}_2^m$ and setting
$F=(-1)^{u\cdot v}$, where the dot product is computed modulo $2$. Then for every
coordinate $i\in\{1,\ldots,2m\}$, the influence of $F$ at $i$ is $1/2$.
\end{theorem}

\begin{theorem}[Fourier regularity without small influences]
\lean{FABL.exists_isFourierRegular_not_hasSmallInfluences_of_even}
\label{FABL.exists_isFourierRegular_not_hasSmallInfluences_of_even}
\leanok
\uses{FABL.BooleanFunction, FABL.BooleanFunction.toReal, FABL.HasSmallInfluences, FABL.IsFourierRegular, FABL.hasSmallInfluences_iff, FABL.influence_innerProductModTwoBoolean_eq_half, FABL.innerProductModTwoBoolean, FABL.isFourierRegular_innerProductModTwoBoolean}
\difficulty{5}
For every positive even integer $n$, there is a Boolean function
$f:\\{-1,1\\}^n\\to\\{\\pm1\\}$ whose nonempty-set Fourier coefficients all have
absolute value at most $2^{-n/2}$. For every real $\\varepsilon<1/2$, this function does
not have all its influences at most $\\varepsilon$.
\end{theorem}

\begin{theorem}[Small stable influences without Fourier regularity]
\lean{FABL.exists_hasSmallStableInfluences_not_isFourierRegular}
\proofsource{Analysis-of-Boolean-Functions-by-Ryan-ODonnell}{pdf-ffbe9c7b3306-p146-b001}
\label{FABL.exists_hasSmallStableInfluences_not_isFourierRegular}
\leanok
\uses{FABL.BooleanFunction, FABL.BooleanFunction.toReal, FABL.HasSmallStableInfluences, FABL.IsFourierRegular, FABL.parityFunction, FABL.parityFunction_not_isFourierRegular_of_lt_one, FABL.stableInfluence_parityFunction}
\difficulty{4}
For every natural number $n>0$ and every $\delta\in[0,1]$, there is a Boolean function
on the $n$-dimensional sign cube whose stable influence at each coordinate, with noise
parameter $1-\delta$, is at most $(1-\delta)^{n-1}$. Moreover, for every real
$\varepsilon<1$, the function has a nonempty Fourier coefficient whose absolute value
exceeds $\varepsilon$.
\end{theorem}

\begin{definition}[Frequency tail]
\lean{FABL.frequencyTail}
\label{FABL.frequencyTail}
\leanok
For a subset $T\subseteq\{0,\ldots,n\}$, define its tail to be the subset of
$\{0,\ldots,n-1\}$ consisting of those $i$ for which $i+1\in T$.
\end{definition}

\begin{theorem}[Recovering a set from its tail frequencies]
\lean{FABL.tailFrequency_frequencyTail_eq}
\label{FABL.tailFrequency_frequencyTail_eq}
\leanok
\uses{FABL.frequencyTail, FABL.tailFrequency}
\difficulty{4}
Let $T$ be a subset of $\\{0,1,\\ldots,n\\}$ that does not contain $0$. Delete the
initial coordinate from the indices in $T$ to obtain a subset of
$\\{0,1,\\ldots,n-1\\}$, then shift those indices up by one; the resulting set is
exactly $T$.
\end{theorem}

\begin{theorem}[Reconstructing a set from its tail]
\lean{FABL.insert_zero_tailFrequency_frequencyTail_eq}
\label{FABL.insert_zero_tailFrequency_frequencyTail_eq}
\leanok
\uses{FABL.frequencyTail, FABL.tailFrequency}
\difficulty{4}
Let $T\subseteq\mathrm{Fin}(n+1)$ be finite and contain $0$. The set obtained by taking
the positive-index elements of $T$, shifting their indices down by one, and adjoining
$0$ is exactly $T$.
\end{theorem}

\begin{theorem}[Vanishing shifted Fourier coefficient]
\lean{FABL.fourierCoeff_leadingCoordinateTimes_tailFrequency_eq_zero}
\label{FABL.fourierCoeff_leadingCoordinateTimes_tailFrequency_eq_zero}
\leanok
\uses{FABL.BooleanFunction, FABL.BooleanFunction.toReal, FABL.firstCoordinateSlice, FABL.fourierCoeff, FABL.fourierCoeff_tailFrequency, FABL.leadingCoordinateTimes, FABL.monomial, FABL.signValue, FABL.tailFrequency}
\difficulty{4}
Let $n$ be a nonnegative integer, let $g$ be a function from the $n$-dimensional Boolean
cube to $\\{−1,1\\}$, and let $S$ be a set of coordinates in $\\{0,\\ldots,n-1\\}$. Form
a function on the $(n+1)$-dimensional cube by multiplying the first coordinate by $g$
evaluated on the remaining coordinates. Its Fourier coefficient at the set obtained by
shifting every coordinate of $S$ up by one is zero.
\end{theorem}

\begin{theorem}[Stable influence and the stability curve]
\lean{FABL.stableInfluence_leadingCoordinateTimes_eq_stabilityCurve}
\label{FABL.stableInfluence_leadingCoordinateTimes_eq_stabilityCurve}
\leanok
\uses{FABL.BooleanFunction, FABL.BooleanFunction.toReal, FABL.card_tailFrequency, FABL.fourierCoeff_leadingCoordinateTimes_insert_zero_tailFrequency, FABL.frequencyTail, FABL.insert_zero_tailFrequency_frequencyTail_eq, FABL.leadingCoordinateTimes, FABL.stabilityCurve, FABL.stableInfluence, FABL.tailFrequency, FABL.zero_notMem_tailFrequency}
\difficulty{6}
Let $n$ be a natural number, let $g:\{−1,1\}^n\to\{−1,1\}$, and let $\rho\in\mathbb R$.
The stable influence at coordinate $0$ of the function on $n+1$ coordinates obtained by
multiplying the leading coordinate by $g$ on the remaining coordinates, after converting
its values to real numbers, equals the stability curve of the real-valued version of $g$
at $\rho$.
\end{theorem}

\begin{theorem}[Majority stable influence and noise stability]
\lean{FABL.stableInfluence_leadingCoordinateTimes_majority_odd_eq_noiseStability}
\label{FABL.stableInfluence_leadingCoordinateTimes_majority_odd_eq_noiseStability}
\leanok
\uses{FABL.BooleanFunction.toReal, FABL.leadingCoordinateTimes, FABL.majority, FABL.noiseStability, FABL.stabilityCurve_eq_noiseStability, FABL.stableInfluence, FABL.stableInfluence_leadingCoordinateTimes_eq_stabilityCurve}
\difficulty{3}
Let $m$ be a natural number and let $0<\delta<1$. The stable influence at parameter
$1-\delta$ and leading coordinate of the function on $2m+2$ coordinates formed by
multiplying the leading coordinate by majority on $2m+1$ coordinates equals the noise
stability of majority on $2m+1$ coordinates at parameter $1-\delta$.
\end{theorem}

\begin{theorem}[Influence bound for leading-coordinate majority]
\lean{FABL.one_sub_sqrt_le_of_hasSmallStableInfluences_leadingCoordinateTimes_majority_odd}
\label{FABL.one_sub_sqrt_le_of_hasSmallStableInfluences_leadingCoordinateTimes_majority_odd}
\leanok
\uses{FABL.BooleanFunction.toReal, FABL.HasSmallStableInfluences, FABL.leadingCoordinateTimes, FABL.majority, FABL.one_sub_sqrt_le_stableInfluence_leadingCoordinateTimes_majority_odd}
\difficulty{3}
Let $m$ be a natural number and let $\varepsilon,\delta\in\mathbb R$ with $0<\delta<1$.
If every coordinate of the real-valued function obtained by multiplying odd majority on
$2m+1$ signs by a new leading coordinate has stable influence at parameter $1-\delta$ at
most $\varepsilon$, then $1-\sqrt{\delta}\leq\varepsilon$.
\end{theorem}

\begin{theorem}[Is balanced majority odd]
\lean{FABL.isBalanced_majority_odd}
\statementsource{Analysis-of-Boolean-Functions-by-Ryan-ODonnell}{
  pdf-ffbe9c7b3306-p045-b003,
  pdf-ffbe9c7b3306-p045-b005
}
\label{FABL.isBalanced_majority_odd}
\leanok
\uses{FABL.BooleanFunction.toReal, FABL.IsBalanced, FABL.fourierCoeff_majority_eq_zero_of_odd_arity_of_even_card, FABL.isBalanced_iff_fourierCoeff_empty_eq_zero, FABL.majority}
\difficulty{3}
% dataset-pipeline: stub -- prose pending build_informal_statements.py
Statement pending regeneration.
\end{theorem}

\begin{theorem}[Regularity after adjoining a coordinate]
\lean{FABL.IsFourierRegular.leadingCoordinateTimes_of_isBalanced}
\label{FABL.IsFourierRegular.leadingCoordinateTimes_of_isBalanced}
\leanok
\uses{FABL.BooleanFunction, FABL.BooleanFunction.toReal, FABL.IsBalanced, FABL.IsFourierRegular, FABL.fourierCoeff_leadingCoordinateTimes_insert_zero_tailFrequency, FABL.fourierCoeff_leadingCoordinateTimes_tailFrequency_eq_zero, FABL.frequencyTail, FABL.insert_zero_tailFrequency_frequencyTail_eq, FABL.isBalanced_iff_fourierCoeff_empty_eq_zero, FABL.leadingCoordinateTimes, FABL.tailFrequency, FABL.tailFrequency_frequencyTail_eq}
\difficulty{5}
Let $n$ be a nonnegative integer and let $g:\\{−1,1\\}^n\\to\\{−1,1\\}$ be a sign-valued
function. Suppose $g$ has mean zero and every nonempty Fourier coefficient of its
real-valued sign representation has absolute value at most $\\varepsilon$, where
$\\varepsilon\\ge 0$. Define a function on $\\{−1,1\\}^{n+1}$ by multiplying the first
input coordinate by the value of $g$ on the remaining coordinates. Every nonempty
Fourier coefficient of this function's real-valued sign representation then has absolute
value at most $\\varepsilon$.
\end{theorem}

\begin{theorem}[Regularity of a Coordinate-Times-Majority Function]
\lean{FABL.isFourierRegular_leadingCoordinateTimes_majority_odd}
\label{FABL.isFourierRegular_leadingCoordinateTimes_majority_odd}
\leanok
\uses{FABL.BooleanFunction.toReal, FABL.IsFourierRegular, FABL.IsFourierRegular.leadingCoordinateTimes_of_isBalanced, FABL.isBalanced_majority_odd, FABL.isFourierRegular_majority_odd, FABL.leadingCoordinateTimes, FABL.majority}
\difficulty{3}
For every natural number $m$, consider the function on $\{\pm1\}^{2m+2}$ that multiplies
the first coordinate by the majority sign of the remaining $2m+1$ coordinates, with ties
assigned the positive sign. Every Fourier coefficient indexed by a nonempty coordinate
set has absolute value at most $1/\sqrt{2m+1}$.
\end{theorem}

\begin{theorem}[A Fourier term is bounded by stable influence]
\lean{FABL.rho_pow_card_sub_one_mul_sq_fourierCoeff_le_stableInfluence}
\proofsource{Analysis-of-Boolean-Functions-by-Ryan-ODonnell}{
  pdf-ffbe9c7b3306-p057-b001,
  pdf-ffbe9c7b3306-p215-b005
}
\label{FABL.rho_pow_card_sub_one_mul_sq_fourierCoeff_le_stableInfluence}
\leanok
\uses{FABL.fourierCoeff, FABL.stableInfluence}
\difficulty{3}
Let $n$ be a natural number, let $f:\{−1,1\}^n\to\mathbb R$, and let $\rho\ge0$. For
every subset $S\subseteq\{1,\ldots,n\}$ and coordinate $i\in S$, the term
$\rho^{|S|-1}\widehat f(S)^2$ is at most the stable influence $\sum_{T\ni
i}\rho^{|T|-1}\widehat f(T)^2$, where $\widehat f(T)$ is the uniform Fourier coefficient
of $f$.
\end{theorem}

\begin{theorem}[Low-degree Fourier regularity]
\lean{FABL.HasSmallStableInfluences.isLowDegreeFourierRegular}
\label{FABL.HasSmallStableInfluences.isLowDegreeFourierRegular}
\leanok
\uses{FABL.HasSmallStableInfluences, FABL.IsLowDegreeFourierRegular, FABL.fourierCoeff, FABL.rho_pow_card_sub_one_mul_sq_fourierCoeff_le_stableInfluence, FABL.stableInfluence}
\difficulty{5}
Let $f:\{-1,1\}^n\to\mathbb R$, let $\epsilon,\delta\in\mathbb R$, and let $k\in\mathbb
N$. Suppose every coordinate has stable influence at correlation parameter $1-\delta$ at
most $\epsilon$, and $0\le\delta<1$. Then every nonempty set $S\subseteq\{1,\ldots,n\}$
with $|S|\le k$ satisfies $|\widehat f(S)|\le\sqrt{\epsilon/(1-\delta)^{\max(k-1,0)}}$.
\end{theorem}

\begin{theorem}[Stable influence at parameter zero]
\lean{FABL.stableInfluence_zero_eq_sq_fourierCoeff_singleton}
\proofsource{Analysis-of-Boolean-Functions-by-Ryan-ODonnell}{pdf-ffbe9c7b3306-p057-b001}
\label{FABL.stableInfluence_zero_eq_sq_fourierCoeff_singleton}
\leanok
\uses{FABL.fourierCoeff, FABL.stableInfluence}
\difficulty{4}
Let $n$ be a natural number, let $f:\\{−1,1\\}^n\\to\\mathbb R$, and let $i$ be a
coordinate. The stable influence of $f$ at parameter $0$ and coordinate $i$ equals the
square of the uniform Fourier coefficient of $f$ on the singleton subset $\\{i\\}$.
\end{theorem}

\begin{theorem}[Unit-noise influences and first-level regularity]
\lean{FABL.hasSmallStableInfluences_one_iff_isLowDegreeFourierRegular_one}
\label{FABL.hasSmallStableInfluences_one_iff_isLowDegreeFourierRegular_one}
\leanok
\uses{FABL.HasSmallStableInfluences, FABL.IsLowDegreeFourierRegular, FABL.fourierCoeff, FABL.stableInfluence_zero_eq_sq_fourierCoeff_singleton}
\difficulty{4}
Let $n$ be a natural number, let $f$ be a real-valued function on the $n$-dimensional
sign cube, and suppose $0\leq\varepsilon$. The stable influences of $f$ at noise
parameter $1$ are all at most $\varepsilon$ if and only if every nonempty subset $S$ of
coordinates with $|S|\leq1$ has Fourier coefficient of absolute value at most
$\sqrt{\varepsilon}$.
\end{theorem}

\begin{theorem}[Fourier coefficient bounded by influence]
\lean{FABL.abs_fourierCoeff_le_influence_of_monotone}
\proofsource{Analysis-of-Boolean-Functions-by-Ryan-ODonnell}{
  pdf-ffbe9c7b3306-p086-b010,
  pdf-ffbe9c7b3306-p048-b002,
  pdf-ffbe9c7b3306-p048-b003
}
\label{FABL.abs_fourierCoeff_le_influence_of_monotone}
\leanok
\uses{FABL.BooleanFunction, FABL.BooleanFunction.toReal, FABL.discreteDerivative, FABL.fourierCoeff, FABL.fourierCoeff_discreteDerivative, FABL.influence, FABL.monomial, FABL.monomial_sq, FABL.sq_discreteDerivative_toReal_eq_self_of_monotone}
\difficulty{4}
Let $n$ be a nonnegative integer and let $f\colon\{\pm1\}^n\to\{\pm1\}$ be monotone in
the coordinatewise order. For any subset $S\subseteq\{1,\ldots,n\}$ and coordinate $i\in
S$, the absolute value of the Fourier coefficient of $f$ on $S$ is at most the influence
of $f$ in coordinate $i$.
\end{theorem}

\begin{theorem}[Is fourier regular and has small influences of monotone]
\lean{FABL.IsLowDegreeFourierRegular.isFourierRegular_and_hasSmallInfluences_of_monotone}
\statementsource{Analysis-of-Boolean-Functions-by-Ryan-ODonnell}{pdf-ffbe9c7b3306-p086-b010}
\label{FABL.IsLowDegreeFourierRegular.isFourierRegular_and_hasSmallInfluences_of_monotone}
\leanok
\uses{FABL.BooleanFunction, FABL.BooleanFunction.toReal, FABL.HasSmallInfluences, FABL.IsFourierRegular, FABL.IsLowDegreeFourierRegular, FABL.abs_fourierCoeff_le_influence_of_monotone, FABL.fourierCoeff, FABL.hasSmallInfluences_iff, FABL.influence, FABL.influence_eq_fourierCoeff_singleton_of_monotone, FABL.influence_nonneg}
\difficulty{4}
% dataset-pipeline: stub -- prose pending build_informal_statements.py
Statement pending regeneration.
\end{theorem}
